% Principal: GDZ PPN236018647, 1917, PDF native 10–11, printed 3–4.
\SourceLeaf{10}\NoSourceHead
\pdfbookmark[1]{Zur Kreistheilung}{kreistheilung}
\HistoricTitle{[ZUR KREISTHEILUNG.]}
\HistoricCitation{Intelligenzblatt der allgemeinen Litteraturzeitung Nr. 66, 1. Juni 1796, S. 554.}
\begin{center}\textbf{III. NEUE ENTDECKUNGEN.}\end{center}

Es ist jedem Anfänger der Geometrie bekannt, dass verschiedene ordentliche Vielecke, namentlich das Dreyeck, Fünfeck, Fünfzehneck, und die, welche durch wiederholte Verdopplung der Seitenzahl eines derselben entstehen, sich geometrisch construiren lassen. So weit war man schon zu \textsc{Euklids} Zeit, und es scheint, man habe sich seitdem allgemein überredet, dass das Gebiet der Elementargeometrie sich nicht weiter erstrecke: wenigstens kenne ich keinen geglückten Versuch, ihre Grenzen auf dieser Seite zu erweitern.

Desto mehr, dünkt mich, verdient die Entdeckung Aufmerksamkeit, dass \textit{ausser jenen ordentlichen Vielecken noch eine Menge anderer, z. B. das Siebzehneck, einer geometrischen Construction fähig ist.} Diese Entdeckung ist eigentlich nur ein Corollarium einer noch nicht ganz vollendeten Theorie von grösserm Umfange, und sie soll, sobald diese ihre Vollendung erhalten hat, dem Publikum vorgelegt werden.

\begin{center}
C. F. \textsc{Gauss}, a. Braunschweig.\\
Stud. der Mathematik zu Göttingen.
\end{center}
\ClosingRule

Es verdient angemerkt zu werden, dass Hr. \textsc{Gauss} jetzt in seinem 18ten Jahre steht, und sich hier in Braunschweig mit eben so glücklichem Erfolge der Philosophie und der classischen Litteratur als der höhern Mathematik gewidmet hat.

\noindent Den 18. April 96
\begin{flushright}E. A. W. \textsc{Zimmermann}, Prof.\end{flushright}
\vfill\begin{flushright}\small 1\textsuperscript{*}\end{flushright}
\clearpage
\SourceLeaf{11}\SourceHead{4}
\begin{center}\small\textbf{BEMERKUNG ZU DER VORSTEHENDEN NOTIZ.}\end{center}
\begingroup\small\par\indent
Diese Ankündigung der Entdeckung der geometrischen Konstruierbarkeit des regelmäßigen Siebzehnecks ist die erste Veröffentlichung von \textsc{Gauss}. \textsc{Sartorius v. Waltershausen} bemerkt\footnote{\textsc{Gauss} zum Gedächtnis, Leipzig 1856, S. 16.}: »Diese Entdeckung, welche er bis zum Ende seines Lebens sehr hoch schätzte, ist es vornehmlich gewesen, welche seinem Leben eine bestimmte Richtung gegebe nhat, denn von jenem Tage an war er fest entschlossen nur der Mathematik sein Leben zu widmen«. Vergl. die erste Eintragung des Tagebuchs vom 30. März 1796. Im ersten Satze des Textes heißt es in der ursprünglichen Veröffentlichung »Viereck« statt »Fünfeck«; vergl. jedoch Werke II, S. 186, Zeile 15 v. u. Im Vorwort zu den \textit{Disqu. Arithm.} (Werke I, S. 6) gedenkt \textsc{Gauss} dieser Anzeige mit den Worten: »Propositum huius operis, ad quod edendum iam annos abhinc quinque publice fidem dederam, . . .«
\begin{flushright}\textsc{Schlesinger}.\end{flushright}
\endgroup
\ClosingRule
\clearpage
